\section{Proofs of Equilibria and Monoculture Results}
\label{app:proofs}
\subsection{Common Cost Function}
\label{app:CommonCost_proof}
\begin{restatable}[Common Cost Functions]{proposition}{CommonCost}
\label{prop:CommonCost}
Suppose each developer $i$ has continuous production function $\Fi$ and cost function $c_i$, where each $c_i:\mathbb{R}_{\ge 0}\to\mathbb{R}_{\ge 0}$ is a strictly increasing bijection with $c_i(0)=0$. Let $c$ be any other such cost function. Then, for each developer $i$, there exists a production function $\tilde{F}_i$ and a bijection between action vectors $\bxi \leftrightarrow\mathbf{z}_i$ such that

$$
c_i\big(\lVert\mathbf{x}_i\rVert_1\big)=c\big(\lVert\mathbf{z}_i\rVert_1\big)
\qquad\text{and}\qquad
\Fi(\mathbf{x}_i)=\tilde{F}_i(\mathbf{z}_i).
$$

Consequently, the game defined by ${(\Fi,c_i)}_{i=1}^{N}$ is equivalent under reparameterization to the game defined by ${(\tilde \Fi,c)}_{i=1}^{N}$: corresponding strategy profiles induce identical payoffs, and their Nash equilibria are in bijective correspondence.
\end{restatable}
\begin{proof}
Define the scalar map

$$
g_i(T) := c_i^{-1}\big(c(T)\big), \qquad T\ge 0.
$$

Since $c_i$ and $c$ are strictly increasing bijections of $\mathbb{R}_{\ge 0}$, so is $g_i$. In particular, $g_i^{-1}$ exists and is itself a strictly increasing bijection of $\mathbb{R}_{\ge 0}$. By construction,
\begin{equation}
c_i\big(g_i(T)\big) = c(T) \qquad \text{for all } T\ge 0. \label{eq:def_c_i}
\end{equation}
Equivalently, substituting $g_i^{-1}(T)$ for $T$ gives
\begin{equation}
c\big(g_i^{-1}(T)\big) = c_i(T) \qquad \text{for all } T\ge 0. \label{eq:def_c}
\end{equation}

For $\mathbf{x}_i \neq \mathbf{0}$ define the \emph{reference-unit} vector

$$
\mathbf{z}_i := \frac{g_i^{-1}\big(\lVert \mathbf{x}_i\rVert_1\big)}{\lVert \mathbf{x}_i\rVert_1}\,\mathbf{x}_i,
$$

and set $\mathbf{z}_i=\mathbf{0}$ when $\mathbf{x}_i=\mathbf{0}$. Thus $\mathbf{x}_i$ is rescaled, keeping its direction fixed, so that $\lVert \mathbf{z}_i\rVert_1 = g_i^{-1}(\lVert \mathbf{x}_i\rVert_1)$. Applying $g_i$ to both sides of this identity gives

$$
g_i\big(\lVert \mathbf{z}_i\rVert_1\big) = \lVert \mathbf{x}_i\rVert_1.
$$

Moreover, by \eqref{eq:def_c},

$$
c\big(\lVert \mathbf{z}_i\rVert_1\big) = c\big(g_i^{-1}(\lVert \mathbf{x}_i\rVert_1)\big) = c_i\big(\lVert \mathbf{x}_i\rVert_1\big),
$$

so cost is exactly preserved.

For $\mathbf{z}\neq\mathbf{0}$, define

$$
\tilde{F}_i(\mathbf{z}) := F_i\!\left(\frac{g_i\big(\lVert \mathbf{z}\rVert_1\big)}{\lVert \mathbf{z}\rVert_1}\,\mathbf{z}\right),
$$

and let $\tilde{F}_{i}(\mathbf{0}) = F_{i}(\mathbf{0})$. 
Evaluating at $\mathbf{z}=\mathbf{z}_i$ for $\mathbf{x}_{i} \neq \mathbf{0}$ using $g_i\big(\lVert \mathbf{z}_i\rVert_1\big) = \lVert \mathbf{x}_i\rVert_1
$ gives,

$$
\tilde{F}_i(\mathbf{z}_i) = F_i\!\left(\frac{\lVert \mathbf{x}_i\rVert_1}{\lVert \mathbf{z}_i\rVert_1}\cdot\frac{g_i^{-1}(\lVert \mathbf{x}_i\rVert_1)}{\lVert \mathbf{x}_i\rVert_1}\,\mathbf{x}_i\right) = F_i(\mathbf{x}_i),
$$

where we used $\lVert \mathbf{z}_i\rVert_1 = g_i^{-1}(\lVert \mathbf{x}_i\rVert_1)$. If $\mathbf{x}_i=\mathbf{0}$, then $\mathbf{z}_i=\mathbf{0}$ and the equality follows directly from $\tilde F_i(\mathbf{0})=F_i(\mathbf{0})$. Thus $F_i(\mathbf{x}_i)=\tilde F_i(\mathbf{z}_i)$ as required.

Finally, the map $\mathbf{x}_i\mapsto\mathbf{z}_i$ is a bijection of $\mathbb{R}_{\ge0}^K$, with inverse

$$
\mathbf{x}_i = \frac{g_i(\lVert \mathbf{z}_i\rVert_1)}{\lVert \mathbf{z}_i\rVert_1}\mathbf{z}_i
$$

for $\mathbf{z}_i\neq\mathbf{0}$, and $\mathbf{x}_i=\mathbf{0}$ when $\mathbf{z}_i=\mathbf{0}$.

Applying this construction to every developer $i$, corresponding strategy profiles $(\mathbf{x}_1,\ldots,\mathbf{x}_N)$ and $(\mathbf{z}_1,\ldots,\mathbf{z}_N)$ induce identical outputs and costs for every developer, and hence identical payoffs. Since the correspondence is a bijection for each developer, best responses are preserved under the reparameterization, and therefore the Nash equilibria of the two games are in bijective correspondence.
\end{proof}


\subsection{Proof of Theorem~\ref{thm:EquilibriumExistence}}
\label{app:EquilibriumExistence_proof}
\EquilibriumExistence*
\begin{proof}
Since $\sigma$ is bounded in $[0,1]$ and $\mathbf{w}\in\Delta^{K-1}$, the total reward of any developer is at most $1$. Moreover, choosing $\mathbf{x}_i=\mathbf{0}$ gives nonnegative utility. Therefore, any action satisfying
\[
c(\lVert\mathbf{x}_i\rVert_1)>1
\]
cannot be a best response. Let $\bar{x}=c^{-1}(1)$. We may therefore restrict each developer's strategy set to
\[
\mathcal{X}_i=\{\mathbf{x}_i\in\mathbb{R}_+^K:\lVert\mathbf{x}_i\rVert_1\leq\bar{x}\},
\]
which is compact and convex.
Since any action outside \(\mathcal X_i\) yields strictly negative utility, while \(\mathbf x_i=\mathbf0\) yields nonnegative utility, no action outside \(\mathcal X_i\) can be a profitable deviation. Thus any Nash equilibrium of the restricted game is also a Nash equilibrium of the original game.
We note that $u_{i}$ is continuous because it is built from the addition and multiplication of continuous functions in $\sigma$, $F_{i}$, and $c$. 

It remains to show that $u_i(\mathbf{x}_i,\mathbf{x}_{-i})$ is  concave in $\mathbf{x}_i$. For $z_{ij,k}=g_{i,k}(\mathbf{x}_i)-g_{j,k}(\mathbf{x}_j)$,
\[
\nabla_{\mathbf{x}_i}^2 \sigma(z_{ij,k})
=
\sigma'(z_{ij,k})\nabla^2 g_{i,k}(\mathbf{x}_i)
+
\sigma''(z_{ij,k})
\nabla g_{i,k}(\mathbf{x}_i)\nabla g_{i,k}(\mathbf{x}_i)^\top.
\]
Since $\sigma'(z)>0$ and
\[
\frac{\sigma''(z)}{\sigma'(z)} \leq \rho_\sigma,
\]
the assumed curvature condition implies
\[
\nabla_{\mathbf{x}_i}^2 \sigma(z_{ij,k})\preceq 0.
\]
Thus every pairwise reward term is concave in $\mathbf{x}_i$. Their nonnegative weighted sum is therefore concave, and since $c$ is convex,
\[
-c(\lVert\mathbf{x}_i\rVert_1)
\]
is concave as well. Hence $u_i(\cdot,\mathbf{x}_{-i})$ is concave (and so quasi-concave) on $\mathcal{X}_i$ for every $\mathbf{x}_{-i}$.

Finally, the utilities are continuous and the strategy sets $\mathcal{X}_i$ are nonempty, compact, and convex. Therefore, by the Debreu--Fan--Glicksberg equilibrium existence theorem, the game admits a pure-strategy Nash equilibrium.
\end{proof}

\subsection{Proof of Lemma~\ref{lem:SymmetricDevelopers}}
\label{app:SymmetricDevelopers_proof}
\SymmetricDevelopers*
\begin{proof}
By Theorem~\ref{thm:EquilibriumExistence}, we may restrict each developer's action space to a compact convex set $\mathcal{X}_i$, and each payoff $u_i$ is continuous and concave in the developer's own action. Hence each best-response correspondence is nonempty, convex-valued, and upper hemicontinuous.

First suppose that $F_i=F_{i'}$. Consider the subset of strategy profiles
\[
\mathcal{X}^{i,i'}
=
\left\{
\bx \in\prod_{j=1}^N\mathcal{X}_j:
\bxi =\bx_{i'}
\right\}.
\]
At any profile $\bx \in\mathcal{X}^{i,i'}$, developers $i$ and $i'$ face identical optimization problems: they have the same production function and cost function, and, because $\bxi=\bx_{i'}$, each faces the same actions and outputs from all other developers. Therefore,
\[
\BR_i(\bxmi)
=
\BR_{i'}(\bx_{-i'}).
\]
Since these best-response sets are identical, nonempty, and convex, we can define a reduced best-response correspondence on $\mathcal{X}^{i,i'}$ that requires developers $i$ and $i'$ to select the same element of their common best-response set, while every other developer selects from its usual best-response set. 
Equivalently, parameterize $\mathcal{X}^{i,i'}$ by a common action $\mathbf y=\mathbf x_i=\mathbf x_{i'}$ and the actions $(\mathbf x_\ell)_{\ell\neq i,i'}$. 
The resulting correspondence is nonempty and convex-valued. 
Moreover, it is upper hemicontinuous because each developer's best-response correspondence is upper hemicontinuous and the restriction to profiles satisfying $\mathbf x_i=\mathbf x_{i'}$ is continuous.
Thus, by Kakutani's fixed-point theorem, there exists a fixed point $\bx^*\in\mathcal{X}^{i,i'}$.
At this fixed point every developer is playing a best response in the original game, and, by construction, $\bxi^*=\bx_{i'}^*$.

Now suppose that $F_i=F$ for every $i$. 
Now consider the fully symmetric subset
\[
\mathcal{X}^{\mathrm{sym}}
=
\left\{
\bx\in\prod_{j=1}^N\mathcal{X}_j:
\bx_1=\cdots=\bx_N
\right\}.
\]
Here, every developer faces the same optimization problem, and therefore faces the same best response problem. 
We can choose the same best response for each developer under this situation. 
As above, the resulting correspondence maps $\mathcal{X}^{\mathrm{sym}}$ into itself and satisfies the conditions of Kakutani's fixed-point theorem.
Therefore, there exists a fixed point with $\bxi$ constant, which implies $\bthetai$ constant, and consequently $H=1$.
\end{proof}

\subsection{Proof of Theorem~\ref{thm:SimilarFHomogeneity}}
\label{app:SimilarFHomogeneity}
\SimilarFHomogeneity*
\begin{proof}
% Establish similarity in pseudogradients
Let $G(\bx)$ be the pseudogradient for the original game with production functions $F_{i}$ and $F_{i'}$ and let $\widetilde{G}(\bx)$ be the pseudogradient for the counterfactual game where $F_{i}$ is replaced with $F_{i'}$ (while everything else stays constant). 

We show that $G$ and $\widetilde{G}$ are close in L2 norm. 
For each developer $\ell$, the corresponding block of the pseudogradient is
\[
G_\ell(\bx)=-\nabla_{\bx_\ell}u_\ell(\bx).
\]
Since
\[
u_\ell(\bx)=\frac{1}{N-1}\sum_{r\neq \ell}\sum_{k=1}^K w_k\sigma\!\left(F_\ell(\bx_\ell)_k-F_r(\bx_r)_k\right)-c(\lVert\bx_\ell\rVert_1),
\]
we have
\[
\nabla_{\bx_\ell}u_\ell(\bx)=\frac{1}{N-1}\sum_{r\neq \ell}\sum_{k=1}^K w_k \sigma' \!\left(F_\ell(\bx_\ell)_k-F_r(\bx_r)_k\right)\nabla_{\bx_\ell}F_\ell(\bx_\ell)_k-\nabla_{\bx_\ell}c(\lVert\bx_\ell\rVert_1).
\]
Since the strategy space is compact and the production functions are continuous, the arguments of $\sigma$ lie in a compact interval.
Since $\sigma$ is twice continuously differentiable, $\sigma'$ and $\sigma''$ are bounded on this interval, so $\sigma'$ is Lipschitz; moreover, since the production functions are continuously differentiable on the compact strategy space, their gradients are uniformly bounded.
Replacing $F_i$ with $F_{i'}$ affects only terms involving developer $i$.
For $\ell=i$, both the argument of $\sigma'$ and the production-function gradient change; for $\ell\neq i$, only the comparison against developer $i$ changes. 
Using $\lVert F_i(\bx_i)-F_{i'}(\bx_i)\rVert_2\leq\epsilon$, $\lVert \nabla F_i(\bx_i)_k-\nabla F_{i'}(\bx_i)_k\rVert_2\leq\epsilon$, and the Lipschitz continuity of $\sigma'$, we can bound these differences term-by-term to obtain
\[
\sup_{\bx\in\mathcal X}\lVert G(\bx)-\widetilde G(\bx)\rVert_2 \leq C\epsilon,
\] 
for some constant $C>0$ independent of $\epsilon$.


% Establish similarity in equilibrium
Next, let $\bx^*$ denote the equilibrium of the original game. 
Because $F_i$ and $F_{i'}$ both satisfy the conditions of Theorem~\ref{thm:EquilibriumExistence}, the counterfactual game also satisfies these conditions. 
By Lemma~\ref{lem:SymmetricDevelopers}, there exists an equilibrium $\widetilde{\bx}^*$ of the counterfactual game such that $\widetilde{\bx}_i^*=\widetilde{\bx}_{i'}^*$. 
We select such an equilibrium $\widetilde{\bx}^*$.
Since each developer's utility is concave and differentiable in its own
strategy, the first-order optimality conditions characterize equilibrium.
That is, because $\bx^*$ and $\widetilde{\bx}^*$ are equilibria of the original and counterfactual games, respectively, we have that 
\[
G(\bx^*)^\top(\widetilde{\bx}^*-\bx^*)\geq0,\qquad \widetilde G(\widetilde{\bx}^*)^\top(\bx^*-\widetilde{\bx}^*)\geq0.
\]
Let $d=\bx^*-\widetilde{\bx}^*$. Adding the two inequalities gives $(G(\bx^*)-\widetilde G(\widetilde{\bx}^*))^\top d\leq0$. 
Recall that $G(\mathbf{x})$ is strongly monotone, and so 
\[
\left(G(\bx)-G(\bx')\right)^\top(\bx-\bx')
\geq
\mu\lVert\bx-\bx'\rVert_2^2.
\]
Setting $\bx=\bx^*$ and $\bx'=\widetilde{\bx}^*$ gives
\[
(G(\bx^*)-G(\widetilde{\bx}^*))^\top d\geq\mu\lVert d\rVert_2^2.
\]
Invoking $(G(\bx^*)-\widetilde G(\widetilde{\bx}^*))^\top d\leq0$ gives that 
\[
\mu\lVert d\rVert_2^2\leq(G(\bx^*)-G(\widetilde{\bx}^*))^\top d\leq(\widetilde G(\widetilde{\bx}^*)-G(\widetilde{\bx}^*))^\top d \leq C\epsilon \lVert d\rVert_2,
\]
where the final inequality follows from Cauchy--Schwarz and the bound on the difference between the pseudogradients.
Rearranging and solving for $d$ gives that 
\[
\lVert\bx^*-\widetilde{\bx}^*\rVert_2\leq\frac{C\epsilon}{\mu}.
\]

Since $\widetilde{\bx}_i^*=\widetilde{\bx}_{i'}^*$,
\[
\lVert\bx_i^*-\bx_{i'}^*\rVert_2\leq \lVert\bx_i^*-\widetilde{\bx}_i^*\rVert_2+\lVert\bx_{i'}^*-\widetilde{\bx}_{i'}^*\rVert_2\leq O(\lVert\bx^*-\widetilde{\bx}^*\rVert_2) =O\left(\frac{\epsilon}{\mu}\right).
\]

% Establish similarity in outputs/monoculture
Finally, to bound $\Hij(\mathbf{x}^*)$, let 
\[
\begin{aligned}
\delta
&=
\lVert\btheta_i^*-\btheta_{i'}^*\rVert_2\\
&=
\lVert F_i(\bx_i^*)-F_{i'}(\bx_{i'}^*)\rVert_2\\
&\leq
\lVert F_i(\bx_i^*)-F_i(\bx_{i'}^*)\rVert_2
+
\lVert F_i(\bx_{i'}^*)-F_{i'}(\bx_{i'}^*)\rVert_2\\
&\leq
L\lVert\bx_i^*-\bx_{i'}^*\rVert_2+\epsilon\\
&=
O\left(\epsilon\left(1+\frac{L}{\mu}\right)\right).
\end{aligned}
\]
Next, recall that
\[
H_{i,i'}(\mathbf{x}^*)=\frac{\langle\btheta_i^*,\btheta_{i'}^*\rangle}{\lVert\btheta_i^*\rVert_2\lVert\btheta_{i'}^*\rVert_2}.
\]
Let $\widehat{\btheta}_i^*=\btheta_i^*/\lVert\btheta_i^*\rVert_2$ and $\widehat{\btheta}_{i'}^*=\btheta_{i'}^*/\lVert\btheta_{i'}^*\rVert_2$. Since both normalized vectors have unit norm,
\[
\lVert\widehat{\btheta}_i^*-\widehat{\btheta}_{i'}^*\rVert_2^2
=2-2\langle\widehat{\btheta}_i^*,\widehat{\btheta}_{i'}^*\rangle
=2(1-H_{i,i'}(\mathbf{x}^*)).
\]
If $\lVert\btheta_i^*\rVert_2,\lVert\btheta_{i'}^*\rVert_2\geq m$, then, adding and subtracting $\btheta_{i'}^*/\lVert\btheta_i^*\rVert_2$ and applying the triangle inequality,
\[
\begin{aligned}
\left\lVert\frac{\btheta_i^*}{\lVert\btheta_i^*\rVert_2}-\frac{\btheta_{i'}^*}{\lVert\btheta_{i'}^*\rVert_2}\right\rVert_2
&=
\left\lVert
\frac{\btheta_i^*-\btheta_{i'}^*}{\lVert\btheta_i^*\rVert_2}
+
\btheta_{i'}^*
\left(
\frac{1}{\lVert\btheta_i^*\rVert_2}
-
\frac{1}{\lVert\btheta_{i'}^*\rVert_2}
\right)
\right\rVert_2 \\
&\leq
\frac{\lVert\btheta_i^*-\btheta_{i'}^*\rVert_2}{\lVert\btheta_i^*\rVert_2}
+
\lVert\btheta_{i'}^*\rVert_2
\left|
\frac{1}{\lVert\btheta_i^*\rVert_2}
-
\frac{1}{\lVert\btheta_{i'}^*\rVert_2}
\right| \\
&=
\frac{\lVert\btheta_i^*-\btheta_{i'}^*\rVert_2}{\lVert\btheta_i^*\rVert_2}
+
\frac{\left|\lVert\btheta_i^*\rVert_2-\lVert\btheta_{i'}^*\rVert_2\right|}{\lVert\btheta_i^*\rVert_2} \\
&\leq
\frac{2\lVert\btheta_i^*-\btheta_{i'}^*\rVert_2}{m}
=
\frac{2\delta}{m}.
\end{aligned}
\]
Therefore,
\[
2(1-H_{i,i'}(\mathbf{x}^*))\leq\frac{4\delta^2}{m^2},
\]
and hence
\[
H_{i,i'}(\mathbf{x}^*)\geq1-\frac{2\delta^2}{m^2}=1-O\left(\frac{\epsilon^2}{m^2}\left(1+\frac{L}{\mu}\right)^2\right).
\]
\end{proof}

\subsection{Proof of Theorem~\ref{thm:SpecializedFDifferentiation}}
\label{app:SpecializedFDifferentiation}
\SpecializedFDifferentiation*

\begin{proof}
Let
\[
    Y_i^*
    :=
    \{\theta_j^*:j\neq i\},
    \qquad
    Y_{i'}^*
    :=
    \{\theta_j^*:j\neq i'\}
\]
denote the rival environments faced by developers $i$ and $i'$ at
equilibrium. Since $\mathbf x^*$ is a Nash equilibrium and $\Theta$
contains all relevant best responses, we have
\[
    \theta_i^*
    =
    \arg\min_{\theta\in\Theta}
    \phi_{i,Y_i^*}(\theta),
    \qquad
    \theta_{i'}^*
    =
    \arg\min_{\theta\in\Theta}
    \phi_{i',Y_{i'}^*}(\theta).
\]
By the assumed interiority of the minimizers,
\[
    \nabla \phi_{i,Y_i^*}(\theta_i^*)=0,
    \qquad
    \nabla \phi_{i',Y_{i'}^*}(\theta_{i'}^*)=0.
\]

We first compare the two developers while holding the rival environment
fixed. Define the counterfactual best response of developer $i'$ to
developer $i$'s rival environment:
\[
    \widetilde\theta_{i'}
    :=
    \arg\min_{\theta\in\Theta}
    \phi_{i',Y_i^*}(\theta).
\]
We will decompose $\lVert \theta^*_i-\theta^*_{i'} \rVert_2$ into $\lVert \theta^*_i-\widetilde{\theta}_{i'} \rVert_2$ and $\lVert \widetilde{\theta}_{i'} -\theta^*_{i'}\rVert_2$.

To bound the first term, we note that $\nabla\phi_{i,Y_i^*}(\theta_i^*)=0$ implies $\nabla C_{i}(\theta_{i}^*) = \nabla V_{Y_{i}^*}(\theta_{i}^*)$. 
As a result, 
\[
\begin{aligned}
    \lVert \nabla\phi_{i',Y_i^*}(\theta_i^*) \rVert_2
    &=
    \lVert\nabla C_{i'}(\theta_i^*)
    -
    \nabla V_{Y_i^*}(\theta_i^*) \rVert_2\\
    &=
    \lVert \nabla C_{i'}(\theta_i^*)
    -
    \nabla C_i(\theta_i^*) \rVert_2 \\
    & \geq \delta_{i,i'} 
\end{aligned}
\]
We additionally note that by interiority, $\nabla\phi_{i',Y_i^*}(\widetilde\theta_{i'})=0$. 
Therefore, combining these two gives that 
\[
 \delta_{i,i'}
    \le
    \left\|
    \nabla\phi_{i',Y_i^*}(\theta_i^*)
    -
    \nabla\phi_{i',Y_i^*}(\widetilde\theta_{i'})
    \right\|_2.
\]
We can upper bound this because $\phi_{i',Y_{i}^*}$ has L-Lipschitz gradient, and so 
\[
 \delta_{i,i'}
    \le
    \left\|
    \nabla\phi_{i',Y_i^*}(\theta_i^*)
    -
    \nabla\phi_{i',Y_i^*}(\widetilde\theta_{i'})
    \right\|_2 \leq  L\|\theta_i^*-\widetilde\theta_{i'}\|_2.
\]
Rearranging gives,
\[
    \|\theta_i^*-\widetilde\theta_{i'}\|_2
    \ge
    \frac{\delta_{i,i'}}{L}.
\]

Next, we need to bound $\lVert\widetilde\theta_{i'}-\theta_{i'}^* \rVert_2$. 
The environments $Y_i^*$ and $Y_{i'}^*$ share the same $N-2$ other rivals and differ only in that $Y_i^*$ contains $\theta_{i'}^*$ whereas $Y_{i'}^*$ contains $\theta_i^*$. 
Thus, for any $\theta\in\Theta$ and coordinate $k$,
\[
\begin{aligned}
&
\left|
\frac{\partial V_{Y_i^*}(\theta)}{\partial\theta_k}
-
\frac{\partial V_{Y_{i'}^*}(\theta)}{\partial\theta_k}
\right| \\
&\qquad =
\frac{w_k}{N-1}
\left|
\sigma'(\theta_k-\theta_{i',k}^*)
-
\sigma'(\theta_k-\theta_{i,k}^*)
\right| \leq \frac{\ell_\sigma w_k}{N-1}
|\theta_{i,k}^*-\theta_{i',k}^*|.
\end{aligned}
\]
Summing across coordinates gives
\[
\begin{aligned}
\left\|
\nabla V_{Y_i^*}(\theta)
-\nabla V_{Y_{i'}^*}(\theta)
\right\|_2^2
&\le
\frac{\ell_\sigma^2}{(N-1)^2}
\sum_{k=1}^K
w_k^2
|\theta_{i,k}^*-\theta_{i',k}^*|^2\\
&\le
\frac{\ell_\sigma^2\|w\|_\infty^2}{(N-1)^2}
\|\theta_i^*-\theta_{i'}^*\|_2^2,
\end{aligned}
\]
and so
\[
    \left\|
    \nabla V_{Y_i^*}(\theta)
    -
    \nabla V_{Y_{i'}^*}(\theta)
    \right\|_2
    \le
    \frac{\ell_\sigma\|w\|_\infty}{N-1}
    \|\theta_i^*-\theta_{i'}^*\|_2.
\]

Because the cost term is the same for both $\phi_{i',Y_{i}^*}$ and $\phi_{i',Y_{i'}^*}$, we have that 
\[
\begin{aligned}
    \left\|
    \nabla\phi_{i',Y_i^*}(\theta_{i'}^*)
    \right\|_2
    &=
    \left\|
    \nabla\phi_{i',Y_i^*}(\theta_{i'}^*)
    -
    \nabla\phi_{i',Y_{i'}^*}(\theta_{i'}^*)
    \right\|_2 \\
    &=
    \left\|
    \nabla V_{Y_{i'}^*}(\theta_{i'}^*)
    -
    \nabla V_{Y_i^*}(\theta_{i'}^*)
    \right\|_2 \\
    &\le
    \frac{\ell_\sigma\|w\|_\infty}{N-1}
    \|\theta_i^*-\theta_{i'}^*\|_2.
\end{aligned}
\]

Since $\phi_{i',Y_i^*}$ is $\mu$-strongly convex and
$\widetilde\theta_{i'}$ is its minimizer, strong monotonicity of its
gradient gives
\[
\begin{aligned}
    \mu
    \|\theta_{i'}^*-\widetilde\theta_{i'}\|_2^2
    &\le
    \left\langle
    \nabla\phi_{i',Y_i^*}(\theta_{i'}^*)
    -
    \nabla\phi_{i',Y_i^*}(\widetilde\theta_{i'}),
    \theta_{i'}^*-\widetilde\theta_{i'}
    \right\rangle \\
    &\le
    \left\|
    \nabla\phi_{i',Y_i^*}(\theta_{i'}^*)
    \right\|_2
    \|\theta_{i'}^*-\widetilde\theta_{i'}\|_2.
\end{aligned}
\]
Therefore,
\[
    \|\theta_{i'}^*-\widetilde\theta_{i'}\|_2
    \le
    \frac{1}{\mu}
    \left\|
    \nabla\phi_{i',Y_i^*}(\theta_{i'}^*)
    \right\|_2 \leq \frac{1}{\mu} \frac{\ell_\sigma\|w\|_\infty}{N-1}
    \|\theta_i^*-\theta_{i'}^*\|_2 = \rho_{N} \lVert \theta_i^*-\theta_{i'}^* \rVert_2,
\]
where
\[
    \rho_N
    =
    \frac{\ell_\sigma\|w\|_\infty}
         {\mu(N-1)}.
\]

Finally, by the triangle inequality,
\[
\begin{aligned}
    \|\theta_i^*-\widetilde\theta_{i'}\|_2
    &\le
    \|\theta_i^*-\theta_{i'}^*\|_2
    +
    \|\theta_{i'}^*-\widetilde\theta_{i'}\|_2 \\
    &\le
    (1+\rho_N)
    \|\theta_i^*-\theta_{i'}^*\|_2.
\end{aligned}
\]
Plugging in the bound for $\lVert \theta^*_i-\widetilde{\theta}_{i'} \rVert_2$ yields
\[
    \|\theta_i^*-\theta_{i'}^*\|_2
    \ge
    \frac{\delta_{i,i'}}
         {L(1+\rho_N)}.
\]

For the cosine-similarity bound, let
\[
    d:=\|\theta_i^*-\theta_{i'}^*\|_2,
    \qquad
    r_i:=\|\theta_i^*\|_2,
    \qquad
    r_{i'}:=\|\theta_{i'}^*\|_2.
\]
For nonzero outputs,
\[
\begin{aligned}
    d^2
    &=
    r_i^2+r_{i'}^2
    -
    2r_ir_{i'}H_{i,i'}(\mathbf x^*)\\
    &=
    (r_i-r_{i'})^2
    +
    2r_ir_{i'}
    \bigl(1-H_{i,i'}(\mathbf x^*)\bigr).
\end{aligned}
\]
Therefore,
\[
    H_{i,i'}(\mathbf x^*)
    =
    1-
    \frac{
        d^2-(r_i-r_{i'})^2
    }{
        2r_ir_{i'}
    }.
\]
Using
\[
    d
    \ge
    \frac{\delta_{i,i'}}
         {L(1+\rho_N)}
\]
and the trivial bound $H_{i,i'}\le1$ gives
\[
H_{i,i'}(\mathbf x^*)
\le
1-
\frac{
\left[
\left(
\frac{\delta_{i,i'}}{L(1+\rho_N)}
\right)^2
-(r_i-r_{i'})^2
\right]_+
}{
2r_ir_{i'}
}.
\]
\end{proof}

\subsection{Proof of Proposition~\ref{prop:DirectionalCharacterization}}
\label{app:DirectionalCharacterization}
\DirectionalCharacterization*
\begin{proof}
First suppose that $D=(\mathbf d_1,\ldots,\mathbf d_N)$ is induced by a pure-strategy Nash equilibrium $\bx^*$ with nonzero outputs.
Let $\alpha_i=\lVert F_i(\bx_i^*)\rVert_2$ and  $\mathbf d_i = \frac{F_i(\bx_i^*)}{\lVert F_i(\bx_i^*)\rVert_2}$, so that $F_i(\bx_i^*)=\alpha_i\mathbf d_i$
for every developer $i$.

For each developer $i$ and capability $k$, define its marginal evaluation weight at
$\bx^*$ as
\[
p_{i,k}^*
=
\frac{w_k}{N-1}
\sum_{j\neq i}
\sigma'\!\left(
F_i(\bx_i^*)_k-F_j(\bx_j^*)_k
\right).
\]
Since $F_i(\bx_i^*)=\alpha_i\mathbf d_i$, this is exactly
\[
p_{i,k}^*
=
p_{i,k}(D,\boldsymbol{\alpha})
=
\frac{w_k}{N-1}
\sum_{j\neq i}
\sigma'\!\left(
\alpha_i d_{i,k}-\alpha_j d_{j,k}
\right).
\]

The gradient of developer $i$'s utility with respect to its investment is therefore
\[
\nabla_{\bx_i}u_i(\bx^*)
=
\nabla F_i(\bx_i^*)^\top
\mathbf p_i(D,\boldsymbol{\alpha})
-
c'(\lVert\bx_i^*\rVert_1)\mathbf 1.
\]
Because $\bx^*$ is an equilibrium, the first-order optimality conditions imply
\[
\nabla F_i(\bx_i^*)^\top
\mathbf p_i(D,\boldsymbol{\alpha})
\leq
c'(\lVert\bx_i^*\rVert_1)\mathbf 1,
\]
with equality in every coordinate $\ell$ such that $x_{i,\ell}^*>0$.
Thus $D$ is locally supportable.

Conversely, suppose that $D$ is locally supportable.
Then there exist $\boldsymbol{\alpha}\in\mathbb R_{++}^N$ and
$\bx=(\bx_1,\ldots,\bx_N)$ such that
\[
F_i(\bx_i)=\alpha_i\mathbf d_i
\]
and
\[
\nabla F_i(\bx_i)^\top
\mathbf p_i(D,\boldsymbol{\alpha})
\leq
c'(\lVert\bx_i\rVert_1)\mathbf 1,
\]
with equality whenever $x_{i,\ell}>0$.

At this profile,
\[
p_{i,k}(D,\boldsymbol{\alpha})
=
\frac{w_k}{N-1}
\sum_{j\neq i}
\sigma'\!\left(
F_i(\bx_i)_k-F_j(\bx_j)_k
\right),
\]
and hence the conditions above are exactly the first-order optimality
conditions for developer $i$'s original payoff $u_i(\bx_i,\bx_{-i})$.

By the curvature argument in the proof of Theorem~\ref{thm:EquilibriumExistence}, \(u_i(\cdot,\bx_{-i})\) is concave on \(\mathbb R_+^K\). 
Therefore, these first-order conditions are sufficient for optimality so $\bx_i$ is a best response to $\bx_{-i}$ for every developer $i$.
It follows that $\bx$ is a pure-strategy Nash equilibrium, and its normalized
output directions are exactly $D$.

Finally, complete monoculture occurs if and only if
\[
\mathbf d_1=\cdots=\mathbf d_N,
\]
so the first characterization follows by setting
$D=(\mathbf d,\ldots,\mathbf d)$.
Similarly, because output vectors are nonnegative, complete specialization
occurs if and only if
\[
\mathbf d_i^\top\mathbf d_j=0
\qquad\forall i\neq j.
\]
The second characterization therefore follows directly.
\end{proof}


\subsection{Proof of Proposition~\ref{prop:AsymptoticBunching}}
\label{app:AsymptoticBunching}
\AsymptoticBunching*

\begin{proof}
Recall that at each time $t$, developer $i$ chooses $\mathbf x_{i,t}$ to maximize the win-rate, minus an unbounded cost function.
The win-rate is at most $1$, while the cost function is increasing and unbounded, so there exists $B<\infty$ such that $c(B)>1$. 
Therefore, at any equilibrium, $\|\mathbf x_{i,t}\|_1\le B$ for every $i,t$. 

Next, we deal with two cases on whether the cumulative investments $\{\mathbf X_i^t\}_{t\ge1}$ are bounded or not, and show that in either case $\|F_i(\mathbf X_i^{t+1})-F_i(\mathbf X_i^t)\|_2 \to 0$.  
First, suppose that $\{\mathbf X_i^t\}_{t\ge1}$ is bounded. 
Because $\mathbf X_i^t$ is the sum of $t$ non-negative numbers, if sum converges, then the individual terms must go to zero; $\sum_t\|\mathbf x_{i,t}\|_1<\infty$, and so $\|\mathbf x_{i,t}\|_1\to0$.
Because $\|\mathbf x_{i,t}\|_1\to0$, then $\lVert \mathbf X_i^{t+1}-\mathbf X_i^t \rVert_1 \to 0$. 
Because $F_i$ is continuously differentiable, its Jacobian is bounded on the compact set containing $\{\mathbf X_i^t\}_{t\ge1}$, and so 
\[
\|F_i(\mathbf X_i^{t+1})-F_i(\mathbf X_i^t)\|_2
\to0.
\]

Instead, suppose that $\{\mathbf X_i^t\}_{t\ge1}$ is unbounded, so  $\|\mathbf X_i^t\|_1\to\infty$.
We know that $\lVert X_i^{t+1}-X_i^{t} \rVert_1 \leq B$.
We can then use the mean-value inequality to show that 
\[
\|F_i(\mathbf X_i^{t+1})-F_i(\mathbf X_i^t)\|_2
\le
B
\sup_{\|\mathbf z\|_1\ge\|\mathbf X_i^t\|_1}
\max_{\ell\in[K]}
\left\|
\frac{\partial F_i(\mathbf z)}{\partial z_\ell}
\right\|_2 .
\]
As $t\to\infty$, $\|\mathbf X_i^t\|_1\to\infty$, and by assumption, we have that $\left\|
\frac{\partial F_i(\mathbf z)}{\partial z_\ell}
\right\|_2 \to 0$.
As a result, $\|F_i(\mathbf X_i^{t+1})-F_i(\mathbf X_i^t)\|_2 \to 0$, so $\|\Delta\boldsymbol{\theta}^t\|_2\to0$.

We next use the catch-up property. 
By Cauchy--Schwarz, we have that 
\[
\beta\|\mathbf g^t\|_2^2
\le
\left\langle
\mathbf g^t,\Delta\boldsymbol{\theta}^t
\right\rangle
\le
\|\mathbf g^t\|_2
\|\Delta\boldsymbol{\theta}^t\|_2.
\]
If $\|\mathbf g^t\|_2>0$, dividing by $\|\mathbf g^t\|_2$ gives
\[
\|\mathbf g^t\|_2
\le
\frac{1}{\beta}
\|\Delta\boldsymbol{\theta}^t\|_2.
\]
Because $\|\Delta\boldsymbol{\theta}^t\|_2 \to 0$ and $\beta$ is a constant, we have that $\lVert \mathbf{g}^t \rVert_2\to 0$. 
By definition of $\mathbf{g}$ this implies that 
\[
g_{i,k}^t
=
\max_{i'}\theta_{i',k}^t-\theta_{i,k}^t
\to0
\qquad\text{for every }i,k.
\]
\end{proof}

\subsection{Proof of Proposition~\ref{prop:PrimitiveCatchup}}
\label{app:PrimitiveCatchup}
\PrimitiveCatchup*
\begin{proof}
Fix a developer $i$, category $k$, and period $t$.  Write
\[
    \theta_{r,k}^t := [F_r(\mathbf X_r^t)]_k,
    \qquad
    \Delta\theta_{r,k}^t
    :=
    \theta_{r,k}^{t+1}-\theta_{r,k}^t,
\]
and let $j\in\arg\max_r \theta_{r,k}^t$ be a developer attaining
the category-$k$ frontier at time $t$. 
That is, $g_{i,k}^t=\theta_{j,k}^t-\theta_{i,k}^t$.
If $g_{i,k}^t=0$, the result follows immediately from coordinatewise
monotonicity of $F_i$, since $\Delta\theta_{i,k}^t\ge 0$.  

Next, suppose $g_{i,k}^t>0$.
We first lower bound the marginal evaluation weight on every category.
For any developers $r,s$, category $q$, and relevant states
$\mathbf x,\mathbf y$, suppose without loss of generality that $[F_r(\mathbf x)]_q \le [F_s(\mathbf y)]_q$.
Then because production functions are $L$-Lipschitz, the gradient is bounded by $L$ and so we have that,
\[
    \mu\bigl([F_s(\mathbf y)]_q-[F_r(\mathbf x)]_q\bigr)
    \le m_{r,q}(\mathbf x)^2 \leq L^2.
\]
Here, the latter term comes from the production functions being $L$-Lipschitz. 
Rearranging gives that 
\[
    \bigl|[F_r(\mathbf x)]_q-[F_s(\mathbf y)]_q\bigr|
    \le \frac{L^2}{\mu}.
\]
Therefore, $|\theta_{i,q}-\theta_{r,q}| \leq L^2/\mu$, and so 
\[
    p_{i,q}
    :=
    \frac{w_q}{N-1}
    \sum_{r\neq i}
    \sigma'\!\left(\theta_{i,q}-\theta_{r,q}\right)
    \ge
    w_{\min}
    \min_{|z|\le L^2/\mu}\sigma'(z)
    =
    \underline p.
\]

Next, let $s_i:=\|\mathbf x_i\|_1$. 
We consider two cases depending on whether or not the trailing developer $i$ surpasses developer $j$ at timestep $t+1$. 
In the first case, suppose $\theta_{i,k}^{t+1}\ge \theta_{j,k}^{t+1}$, then 
\[
    \Delta\theta_{i,k}^t
    =
    \theta_{i,k}^{t+1}-\theta_{i,k}^t
    \ge
    \theta_{j,k}^{t+1}-\theta_{i,k}^t \geq \theta_{j,k}^{t}-\theta_{i,k}^t
    =
    g_{i,k}^t.
\]
Here, the last inequality comes from the fact that investments are non-negative and production functions are non-decreasing. 
As a result, we have that $ \Delta\theta_{i,k}^t\ge \beta g_{i,k}^t$ for any $\beta \leq 1$. 

For the second case, suppose that $\theta_{i,k}^{t+1}<\theta_{j,k}^{t+1}$.
Then we can bound this difference by noting that $\mu (\theta_{j,k}^{t+1}-\theta_{i,k}^{t+1}) \leq m_{i,k}(\mathbf{x})^2$. 
We bound $m_{i,k}(\mathbf{x})$ by noting that at equilibrium, we have that $c'(s_i)=\sum_{q=1}^K p_{i,q}\frac{\partial [F_i(\mathbf X_i^{t+1})]_q} {\partial x_\ell}$, and so 
\[
    c'(s_i)
    \ge
    p_{i,k}
    \frac{\partial [F_i(\mathbf X_i^{t+1})]_k}
         {\partial x_\ell}
    \ge
    \underline p\,
    m_{i,k}(\mathbf X_i^{t+1}).
\]
We note that $\theta_{j,k}^{t+1}-\theta_{i,k}^{t+1} > 0$ and so $ m_{i,k}(\mathbf X_i^{t+1})^2 \geq \mu (\theta^{t+1}_{j,k}-\theta^{t+1}_{i,k}) > 0$, and so zero investment cannot satisfy the first-order optimality condition; hence $s_i>0$ and such an
$\ell$ exists.
We next bound $c'(s_i)$ as 
\[
    c'(s_i)
    =
    \int_0^{s_i}c''(u)\,du
    \le Ms_i.
\]
Therefore
\[
Ms_{i} \geq c'(s_{i}) \geq \underline{p} m_{i,k}(\mathbf{X}_{i,}^{t+1}).
\]
Rearranging gives 
\[
    s_i
    \ge
    \frac{\underline p}{M}
    m_{i,k}(\mathbf X_i^{t+1}).
\]
We finally substitute this back in and simplify with $\Delta \theta^{t}_{i,k}$ using the concavity of $F_{i}$ to get 
\[
\begin{aligned}
    \Delta\theta_{i,k}^t
    &=
    [F_i(\mathbf X_i^{t+1})]_k
    -
    [F_i(\mathbf X_i^t)]_k\\
    &\ge
    \nabla [F_i(\mathbf X_i^{t+1})]_k^\top
    \mathbf{x}_i \\
    &\ge
    m_{i,k}(\mathbf X_i^{t+1})\,s_i \\ 
    &\ge 
    \frac{\underline p}{M}
    m_{i,k}(\mathbf X_i^{t+1})^2 \\ 
    &\ge 
    \frac{\mu \underline p}{M} (\theta_{j,k}^{t+1}-\theta_{i,k}^{t+1}) \\ 
    &\ge 
    \frac{\mu \underline p}{M} (g_{i,k}^t-\Delta \theta_{i,k}^t).
\end{aligned}
\]
Here, the last inequality comes from $\theta_{j,k}^{t+1}-\theta_{i,k}^{t+1} = 
    g_{i,k}^t
      +\Delta\theta_{j,k}^t
      -\Delta\theta_{i,k}^t \ge
    g_{i,k}^t-\Delta\theta_{i,k}^t$.
Rearranging gives 
\[
    \Delta\theta_{i,k}^t
    \ge
    \frac{\underline p\mu}
         {M+\underline p\mu}
    g_{i,k}^t
    =
    \beta g_{i,k}^t.
\]
\end{proof}
